% F. J. Marshall - J. R. Capablanca, Match, 1909, with Capablanca's notes.
% Build with: python3 ../../tools/chessbook.py game-01-marshall-capablanca.tex
\documentclass[10pt,twoside]{article}
\usepackage[paperwidth=487.56bp,paperheight=691.65bp,top=72pt,bottom=44pt,inner=50pt,outer=48pt,
            headheight=13pt,headsep=20pt,footskip=0pt]{geometry}
\usepackage[T1]{fontenc}
\usepackage[lining]{ebgaramond}
\usepackage{graphicx}
\usepackage{microtype}
\usepackage{tikz}
\usepackage{fancyhdr}
\usepackage[hidelinks,pdftitle={Marshall - Capablanca, Match, 1909},
            pdfauthor={J. R. Capablanca}]{hyperref}

\newcommand{\textsizebook}{\fontsize{10.5}{12.6}\selectfont}
\AtBeginDocument{\textsizebook}
\setlength{\columnsep}{11pt}
\setlength{\parindent}{9pt}
\setlength{\parskip}{0pt}
\setlength{\emergencystretch}{1.5em}
\raggedbottom

\newlength{\boardwidth}
\setlength{\boardwidth}{160pt}

% A main line move, as a bold heading above its note
\newcommand{\Move}[1]{\par\addvspace{8pt}\noindent\phantomsection{\bfseries #1}\par\nopagebreak}
% A diagram in the text
\newcommand{\ChessDiagram}[1]{\par\addvspace{7pt}{\centering #1\par}\addvspace{5pt}}

\pagestyle{fancy}
\fancyhf{}
\renewcommand{\headrulewidth}{0pt}
\fancyhead[LE,RO]{\thepage}
\fancyhead[CE]{Queen's Gambit Declined}
\fancyhead[CO]{Marshall -- Capablanca}

\game{d4 d5 c4 e6 Nc3 Nf6 Bg5 Be7 e3 Ne4 Bxe7 Qxe7 Bd3 Nxc3 bxc3 Nd7 Nf3 O-O cxd5 exd5 Qb3 Nf6
      a4 c5 Qa3 b6 a5 Bb7 O-O Qc7 Rfb1 Nd7 Bf5 Rfc8 Bxd7 Qxd7 a6 Bc6 dxc5 bxc5 Qxc5 Rab8 Rxb8
      Rxb8 Ne5 Qf5 f4 Rb6 Qxb6}

\begin{document}

\twocolumn[{%
\thispagestyle{empty}
\begin{tikzpicture}[remember picture,overlay]
  \node[circle,draw=black,line width=1.6pt,fill=black!20,minimum size=112pt] (seal)
       at ([xshift=75pt,yshift=-75pt]current page.north west) {};
  \draw[line width=0.9pt] ([shift={(-48pt,-30pt)}]seal.center) -- ([shift={(44pt,34pt)}]seal.center);
  \node[rotate=37] at ([shift={(-11pt,14pt)}]seal.center) {\fontsize{25}{25}\selectfont\bfseries Game};
  \node[rotate=37] at ([shift={(10pt,-16pt)}]seal.center) {\fontsize{38}{38}\selectfont\bfseries 1};
  \node[anchor=north east] at ([xshift=-14pt,yshift=-24pt]current page.north east)
       {\board[after=4...Be7,width=108pt,bare]};
\end{tikzpicture}

\vspace*{38pt}
\begin{center}
  \begin{minipage}{235pt}\centering
    {\fontsize{22}{25}\selectfont\bfseries Queen's Gambit Declined\par}
    \vspace{4pt}\rule{\linewidth}{1pt}\\[-9.5pt]\rule{\linewidth}{0.4pt}
  \end{minipage}

  \vspace{14pt}
  {\fontsize{19}{22}\selectfont\bfseries Marshall -- Capablanca\par}
  \vspace{10pt}
  {\fontsize{15}{18}\selectfont\bfseries Match, 1909\par}
  \vspace{8pt}
  {\fontsize{11}{13}\selectfont White: F. J. Marshall \quad Black: J. R. Capablanca\par}
  \vspace{10pt}
  {\centering\fontsize{8.5}{10}\selectfont\itshape
   Notes by J. R. Capablanca, from \textup{Chess Fundamentals} (1921), which is in the public domain.\\
   Descriptive notation has been converted to algebraic. Page references are to the book.\\
   Click any move or diagram to open it on Lichess.\par}
\end{center}
\vspace{10pt}
}]
\textsizebook

\mainline{1.d4 d5 2.c4 e6 3.Nc3 Nf6 4.Bg5 Be7 5.e3 Ne4}

I had played this defence twice before in the match with good results, and although I lost this game I still played it until the very last game, when I changed my tactics. The reason was my total lack of knowledge of the different variations in this opening, coupled with the fact that I knew that Dr. E. Lasker had been successful with it against Marshall himself in 1907. I thought that since Dr. Lasker had played it so often, it should be good. The object is to exchange a couple of pieces and at the same time to bring about a position full of possibilities and with promising chances of success once the end-game stage is reached. On general principles it should be wrong, because the same Knight is moved three times in the opening, although it involves the exchange of two pieces. In reality the difficulty in this variation, as well as in nearly all the variations of the Queen's gambit, lies in the slow development of Black's Queen Bishop. However, whether this variation can or cannot be safely played is a question still to be decided, and it is outside the scope of this book. I may add that at present my preference is for a different system of development, but it is not unlikely that I should some time come back to this variation.

\mainline{6.Bxe7 Qxe7 7.Bd3}

\moves{7.cxd5} is preferable for reasons that we shall soon see.

\mainline{7...Nxc3 8.bxc3 Nd7}

Now \moves{8...dxc4} would be a better way to develop the game. The idea is that after \moves[as=a]{8...dxc4 9.Bxc4 b6}, followed by \figB b7, would give Black's Bishop a powerful range. For this variation see the eleventh game of the match.

\mainline{9.Nf3 O-O}

No longer would \moves[as=a]{9...dxc4 10.Bxc4 b6} be good, because \moves[from=a]{11.Bb5} would prevent \figB b7 on account of \figN e5.

\mainline{10.cxd5 exd5 11.Qb3 Nf6 12.a4 c5}

Played with the intention of obtaining the majority of Pawns on the Queen's side. Yet it is doubtful whether this move is good, since it leaves Black's Queen's-side Pawns disrupted in a way. The safer course would have been to play \moves{12...c6}.

\mainline{13.Qa3 b6}

\board

This exposes Black to further attack by \moves{14.a5} without any compensation for it. If I had to play this position nowadays I would simply play \moves[as=a]{13...Re8}. Then after \moves[from=a]{14.Qxc5 Qxc5} would follow, and I believe that Black would regain the Pawn. If, instead, White played \moves[from=a,as=b]{14.dxc5} then \moves[from=b]{14...Bg4} would give Black an excellent game.

\mainline{14.a5 Bb7 15.O-O Qc7 16.Rfb1 Nd7}

\board

Black's position was bad and perhaps lost in any case, but the text move makes matters worse. As a matter of fact I never saw White's reply \moves{17.Bf5}. It never even passed through my mind that this was threatened. Black's best move would have been \moves{16...Rfb8}. If that loses, then any other move would lose as well.

\mainline{17.Bf5 Rfc8}

From bad to worse. \moves{17...Nf6} offered the only hope.

\mainline{18.Bxd7 Qxd7 19.a6 Bc6 20.dxc5 bxc5 21.Qxc5 Rab8}

The game was lost. One move was as good as another.

\mainline{22.Rxb8 Rxb8 23.Ne5 Qf5 24.f4 Rb6 25.Qxb6!}

\noindent\textbf{Resigns.}

Of course, if \moves{25.Nxc6 Rb1+} would have drawn. The text move is pretty and finishes quickly. A well-played game on Marshall's part.

\end{document}
